\documentclass[10pt,a4paper]{article}

\usepackage[T1]{fontenc}
\usepackage[utf8]{inputenc}
\usepackage[ngerman]{babel}
\usepackage[a4paper,margin=14mm]{geometry}
\usepackage{lmodern}
\usepackage{xcolor}
\usepackage{tabularx}
\usepackage{array}
\usepackage{listings}
\usepackage{microtype}

\pagestyle{empty}
\setlength{\parindent}{0pt}
\setlength{\parskip}{0pt}
\renewcommand{\familydefault}{\sfdefault}

\definecolor{Navy}{HTML}{17324D}
\definecolor{Blue}{HTML}{3B6EA5}
\definecolor{Sky}{HTML}{EAF3FA}
\definecolor{Mint}{HTML}{E8F5EF}
\definecolor{Gold}{HTML}{F3B33D}
\definecolor{Ink}{HTML}{1E2935}
\definecolor{Muted}{HTML}{617080}
\definecolor{Line}{HTML}{D8E1E8}
\definecolor{CodeBg}{HTML}{F5F7F9}

\color{Ink}

\newcommand{\sectiontitle}[2]{%
  \vspace{2.3mm}
  {\color{#1}\bfseries\large #2}\par
  \vspace{0.8mm}
  {\color{#1}\rule{\linewidth}{0.7pt}}\par
  \vspace{1.2mm}
}

\newcommand{\solutionrow}[3]{%
  \textcolor{Blue}{\bfseries #1} &
  \ttfamily\small #2 &
  \small #3 \\[1.8mm]
}

\lstset{
  basicstyle=\ttfamily\small,
  backgroundcolor=\color{CodeBg},
  frame=single,
  rulecolor=\color{Line},
  framesep=2.5mm,
  xleftmargin=1mm,
  xrightmargin=1mm,
  aboveskip=1mm,
  belowskip=1mm,
  showstringspaces=false,
  columns=fullflexible
}

\begin{document}

\colorbox{Navy}{%
  \parbox{\dimexpr\linewidth-2\fboxsep\relax}{%
    \vspace{3mm}
    \color{white}
    \begin{tabularx}{\linewidth}{@{}Xr@{}}
      {\small\bfseries INFORMATIK · KLASSE 9} &
      \colorbox{Gold}{\color{Navy}\small\bfseries\strut LÖSUNGEN}
    \end{tabularx}
    \vspace{2.5mm}

    {\fontsize{22}{25}\selectfont\bfseries Sicherungsblatt: Aufgabe 2}\par
    \vspace{1mm}
    {\large Klassen analysieren und verändern}
    \vspace{3mm}
  }%
}

\vspace{3mm}
\colorbox{Sky}{%
  \parbox{\dimexpr\linewidth-2\fboxsep\relax}{%
    \vspace{1.8mm}
    \textcolor{Blue}{\bfseries Ziel:}
    Du kannst Objekte einer Klasse erzeugen, Methoden aufrufen und aus
    Java-Code eine Klassenkarte ableiten.
    \vspace{1.8mm}
  }%
}

\sectiontitle{Blue}{1 · Musterlösung zu a): Programmausgabe}

\begin{tabularx}{\linewidth}{@{}X X@{}}
  \small Die Anweisungen werden von oben nach unten ausgeführt. Zuerst werden
  zwei Hund-Objekte erzeugt, danach werden Methoden aufgerufen. &
  \colorbox{CodeBg}{\parbox{\dimexpr\linewidth-4mm\relax}{%
    \ttfamily\small
    Der Hund heißt Wuffti und ist 5 Jahre alt.\par
    Der Hund heißt Schnuffi und ist 8 Jahre alt.\par
    Schnuffi: Wuff wuff
  }}
\end{tabularx}

\sectiontitle{Blue}{2 · Musterlösung zu b): Programm.java erklären}

\renewcommand{\arraystretch}{1.03}
\begin{tabularx}{\linewidth}{@{}>{\centering\arraybackslash}p{7mm} >{\raggedright\arraybackslash}p{77mm} X@{}}
  \solutionrow{1}{Hund petersHund = new Hund(5, \textquotedbl{}Wuffti\textquotedbl{});}{Erzeugt einen fünf Jahre alten Hund namens Wuffti und speichert ihn in \texttt{petersHund}.}
  \solutionrow{2}{Hund inasHund = new Hund(8, \textquotedbl{}Schnuffi\textquotedbl{});}{Erzeugt ein zweites Hund-Objekt mit Alter 8 und Name Schnuffi. Es wird in \texttt{inasHund} gespeichert.}
  \solutionrow{3}{petersHund.zeigeDaten();}{Ruft \texttt{zeigeDaten()} für Peters Hund auf und gibt seine Daten aus.}
  \solutionrow{4}{inasHund.zeigeDaten();}{Ruft \texttt{zeigeDaten()} für Inas Hund auf und gibt ihre Daten aus.}
  \solutionrow{5}{inasHund.belle();}{Ruft \texttt{belle()} für Inas Hund auf und gibt den Belltext mit seinem Namen aus.}
\end{tabularx}

\sectiontitle{Blue}{3 · Musterlösung zu c): Pauls Hund ergänzen}

\begin{tabularx}{\linewidth}{@{}p{93mm}X@{}}
  \colorbox{CodeBg}{\parbox{87mm}{%
    \vspace{1.5mm}
    \ttfamily\small
    Hund paulsHund = new Hund(2, \textquotedbl{}Spike\textquotedbl{});\par
    paulsHund.belle();\par
    paulsHund.zeigeDaten();\par
    \vspace{1.5mm}
  }} &
  \small
  \textcolor{Blue}{\bfseries Wichtig:}
  Alle drei Zeilen gehören in \texttt{Programm.java}. Das Objekt wird zuerst
  erzeugt. Erst danach können seine Methoden aufgerufen werden.
\end{tabularx}

\sectiontitle{Blue}{4 · Musterlösung zu d): Klassenkarte Hund}

\begin{tabularx}{\linewidth}{@{}p{76mm}X@{}}
  \fcolorbox{Blue}{white}{%
    \parbox{69mm}{%
      \centering
      \vspace{1mm}
      {\large\bfseries Hund}\par
      \vspace{1mm}
      {\color{Blue}\rule{\linewidth}{0.6pt}}\par
      \vspace{1mm}
      \raggedright
      \texttt{alter: int}\par
      \texttt{name: String}\par
      \vspace{1mm}
      {\color{Blue}\rule{\linewidth}{0.6pt}}\par
      \vspace{1mm}
      \texttt{Hund(par1: int, par2: String)}\par
      \texttt{zeigeDaten(): void}\par
      \texttt{belle(): void}\par
      \vspace{1mm}
    }%
  } &
  \small
  \textcolor{Blue}{\bfseries Aufbau der Klassenkarte}

  \vspace{1mm}
  \textbf{oben:} Klassenname

  \textbf{mittig:} Attribute mit Datentyp

  \textbf{unten:} Konstruktor und Methoden

  \vspace{1mm}
  \texttt{void} bedeutet, dass die Methode keinen Wert zurückgibt.
\end{tabularx}

\vspace{2.5mm}
\colorbox{Mint}{%
  \parbox{\dimexpr\linewidth-2\fboxsep\relax}{%
    \vspace{1.4mm}
    \textcolor{Blue}{\bfseries Merksatz:}
    Die Klasse \texttt{Hund} ist der Bauplan. \texttt{petersHund},
    \texttt{inasHund} und \texttt{paulsHund} sind verschiedene Objekte
    dieser Klasse. Jedes Objekt besitzt eigene Attributwerte.
    \vspace{1.4mm}
  }%
}

\vfill
{\color{Line}\rule{\linewidth}{0.5pt}}\par
\vspace{1.5mm}
\begin{tabularx}{\linewidth}{@{}Xr@{}}
  {\small\color{Muted}Klasse 9 · Modellierung und Programmierung} &
  {\small\color{Muted}Sicherungsblatt · Aufgabe 2}
\end{tabularx}

\end{document}
